\section{Label Overlap Stress Test -- 2D DPTable}

\subsection{Scenario A: Same-column annotations with labels}

\begin{animation}[id="overlap-2d-col", label="Same-column annotations"]
\shape{dp}{DPTable}{rows=5, cols=5, label="$dp[i][j]$"}

\step
\apply{dp.cell[0][2]}{value=3}
\apply{dp.cell[1][2]}{value=7}
\apply{dp.cell[2][2]}{value=12}
\apply{dp.cell[3][2]}{value=18}
\apply{dp.cell[4][2]}{value=25}
\narrate{Column 2 filled. Now add vertical annotations along the same column.}

\step
\annotate{dp.cell[1][2]}{label="dp[0][2]+4=7", arrow_from="dp.cell[0][2]", color=info}
\annotate{dp.cell[2][2]}{label="dp[1][2]+5=12", arrow_from="dp.cell[1][2]", color=info}
\annotate{dp.cell[3][2]}{label="dp[2][2]+6=18", arrow_from="dp.cell[2][2]", color=info}
\annotate{dp.cell[4][2]}{label="dp[3][2]+7=25", arrow_from="dp.cell[3][2]", color=info}
\narrate{Four vertical annotations stacking in column 2 -- labels fight for space.}
\end{animation}

\subsection{Scenario B: Same-row annotations with labels}

\begin{animation}[id="overlap-2d-row", label="Same-row annotations"]
\shape{dp2}{DPTable}{rows=5, cols=5, label="$dp[i][j]$"}

\step
\apply{dp2.cell[2][0]}{value=1}
\apply{dp2.cell[2][1]}{value=4}
\apply{dp2.cell[2][2]}{value=8}
\apply{dp2.cell[2][3]}{value=13}
\apply{dp2.cell[2][4]}{value=19}
\narrate{Row 2 filled. Now add horizontal annotations along the same row.}

\step
\annotate{dp2.cell[2][1]}{label="dp[2][0]+3=4", arrow_from="dp2.cell[2][0]", color=good}
\annotate{dp2.cell[2][2]}{label="dp[2][1]+4=8", arrow_from="dp2.cell[2][1]", color=good}
\annotate{dp2.cell[2][3]}{label="dp[2][2]+5=13", arrow_from="dp2.cell[2][2]", color=good}
\annotate{dp2.cell[2][4]}{label="dp[2][3]+6=19", arrow_from="dp2.cell[2][3]", color=good}
\narrate{Four horizontal annotations along row 2 -- labels overlap horizontally.}
\end{animation}

\subsection{Scenario C: Diagonal annotations crossing each other}

\begin{animation}[id="overlap-2d-diag", label="Crossing diagonal annotations"]
\shape{dp3}{DPTable}{rows=5, cols=5, label="$dp[i][j]$"}

\step
\apply{dp3.cell[0][0]}{value=0}
\apply{dp3.cell[1][1]}{value=5}
\apply{dp3.cell[2][2]}{value=12}
\apply{dp3.cell[3][3]}{value=20}
\apply{dp3.cell[4][4]}{value=30}
\apply{dp3.cell[0][4]}{value=2}
\apply{dp3.cell[1][3]}{value=8}
\apply{dp3.cell[3][1]}{value=15}
\apply{dp3.cell[4][0]}{value=22}
\narrate{Diagonal cells filled. Now add crossing diagonal annotations.}

\step
\annotate{dp3.cell[2][2]}{label="dp[0][0]+12=12", arrow_from="dp3.cell[0][0]", color=info}
\annotate{dp3.cell[2][2]}{label="dp[1][1]+7=12", arrow_from="dp3.cell[1][1]", color=good}
\annotate{dp3.cell[2][2]}{label="dp[0][4]+10=12", arrow_from="dp3.cell[0][4]", color=warn}
\annotate{dp3.cell[2][2]}{label="dp[4][0]+val=12", arrow_from="dp3.cell[4][0]", color=error}
\narrate{Four diagonal arrows all converging on cell [2][2] from different corners.}

\step
\annotate{dp3.cell[4][4]}{label="dp[0][0]+30=30", arrow_from="dp3.cell[0][0]", color=path}
\annotate{dp3.cell[0][4]}{label="dp[4][0]+val=2", arrow_from="dp3.cell[4][0]", color=muted}
\narrate{Long diagonal arrows crossing the entire grid in opposite directions.}
\end{animation}

\subsection{Scenario D: Four annotations pointing to dp[2][2] from 4 directions}

\begin{animation}[id="overlap-2d-4dir", label="4 directions into one cell"]
\shape{dp4}{DPTable}{rows=5, cols=5, label="$dp[i][j]$"}

\step
\apply{dp4.cell[2][2]}{value=20}
\apply{dp4.cell[2][1]}{value=15}
\apply{dp4.cell[2][3]}{value=17}
\apply{dp4.cell[1][2]}{value=13}
\apply{dp4.cell[3][2]}{value=16}
\narrate{Center cell and its four orthogonal neighbors filled.}

\step
\annotate{dp4.cell[2][2]}{label="dp[2][1]+5=20", arrow_from="dp4.cell[2][1]", color=info}
\annotate{dp4.cell[2][2]}{label="dp[2][3]+3=20", arrow_from="dp4.cell[2][3]", color=good}
\annotate{dp4.cell[2][2]}{label="dp[1][2]+7=20", arrow_from="dp4.cell[1][2]", color=warn}
\annotate{dp4.cell[2][2]}{label="dp[3][2]+4=20", arrow_from="dp4.cell[3][2]", color=error}
\narrate{Four annotations from left, right, above, and below all targeting cell [2][2]. Labels collide at the center.}
\end{animation}
